\begin{onehalfspace}

La risposta al rumore può agire sull'hardware, sulla stima dell'output, sulla
codifica dell'informazione o sul circuito. Questa sezione confronta tali leve e
isola quelle effettivamente valutate; la rassegna estesa è nella
Sezione~\ref{app:strategie_estese} dell'Appendice~\ref{app:teoria}.

\section{Le Quattro Strategie: Panoramica e Scelta}

La \textbf{soppressione} riduce l'interazione con l'ambiente. La \textbf{mitigazione} usa campioni aggiuntivi per stimare l'output ideale dopo il calcolo. La \textbf{fault tolerance} protegge i qubit logici usando molti qubit fisici: sotto soglia, aumentare la distanza del codice riduce l'errore~\cite{aharonov1997fault}. I \textbf{metodi appresi} svolgono invece compiti classici, come classificare gli output o decodificare le sindromi. Il machine learning può quindi assistere la \ac{QEC}, ma non realizza da solo la protezione dello stato.

\begin{table}[H]
\centering
\small
\begin{tabular}{@{}llll@{}}
\toprule
\textbf{Strategia} & \textbf{Quando agisce} & \textbf{Costo dominante} & \textbf{Uso qui} \\
\midrule
Soppressione & prima/durante il calcolo & controllo hardware & no \\
Mitigazione & dopo il calcolo & campioni aggiuntivi & confronto ZNE \\
Fault tolerance/QEC & durante il calcolo & qubit e cicli & sì, simulata \\
Metodi appresi & post-processing/decoder & dati e calcolo classico & valutati \\
\bottomrule
\end{tabular}
\caption{Confronto delle strategie anti-rumore nel perimetro della tesi.}
\label{tab:strategie_confronto}
\end{table}

L'utilità del modello appreso viene verificata in entrambe le fasi. Nella prima, il confronto con TOP-1 e TOP-4 usa lo stesso istogramma. Nella seconda, la rete interviene nella decodifica delle sindromi. La \ac{ZNE} fornisce un confronto aggiuntivo.

\section{Machine Learning Classico per la Robustezza al Rumore}
\label{sec:qml_overview}

Random Forest, \ac{SVM} e \ac{MLP} sono qui modelli \textbf{classici applicati
a dati quantistici}, non \ac{QML}: addestramento e inferenza non impiegano un
circuito quantistico. Possibili usi comprendono caratterizzazione del rumore,
ottimizzazione circuitale e decodifica delle sindromi; il lavoro verifica il
post-processing dell'istogramma e, separatamente, un correttore residuo del
decoder. Le altre possibilità sono documentate nell'Appendice A.

\section[Correzione degli Errori Quantistici e Ruolo del Machine Learning]{Correzione degli Errori Quantistici e Ruolo del Machine Learning}
\label{sec:qec_overview}

La \ac{QEC} distribuisce un qubit logico su più qubit fisici e ricava
informazione sull'errore senza misurare lo stato logico. Il ciclo essenziale è:
\begin{enumerate}
    \item \textbf{codifica} dello stato $\ket{\psi_L}$;
    \item \textbf{misura degli stabilizzatori}, che produce una sindrome;
    \item \textbf{decodifica e recupero}, analitici o assistiti da un modello.
\end{enumerate}
Il codice di Shor $[[9,1,3]]$ dimostra la correzione di un errore arbitrario;
Steane $[[7,1,3]]$ riduce l'overhead con una costruzione \ac{CSS}; il surface
code usa una griglia locale e sopprime l'errore logico sotto una soglia che
dipende da circuito, rumore e decoder. Stati logici e proprietà complete sono
riportati nella Sezione~\ref{app:qec_richiami_teoria}.

\subsection{Machine Learning come Possibile Componente del Decoder QEC}
\label{sec:qec_ml}

La sindrome conserva informazione strutturata sull'errore e rende sensato un
confronto fra decoder analitico e appreso sugli stessi campioni
\cite{torlaiMelko2017,alphaqubit2024}. Nella
Sezione~\ref{sec:decodifica_appresa} il modello non sostituisce con vantaggio il
matching; come correttore residuo, invece, recupera parte delle correlazioni che
il grafo analitico non rappresenta. Il risultato resta circoscritto al modello
simulato.

\section{Motivazione Pratica: il Contributo di questa Tesi}
\label{sec:motivazione_ml}

Nella prima fase, M1 prova l'esito più frequente e TOP-4 prova i primi quattro. M2 usa un classificatore per decidere se attivare la stessa TOP-4. Confrontare M2 con TOP-4 senza filtro chiarisce l'origine del miglioramento: cercare più candidati è utile, mentre il classificatore non aggiunge un beneficio. Questo risultato motiva lo studio della correzione durante il calcolo.

\section{Una Strategia sul Circuito: la Trasformata di Fourier Approssimata}
\label{sec:aqft_strategia}

Una leva complementare riduce il circuito prima che agisca il rumore. La
\ac{AQFT} elimina le rotazioni controllate più piccole: introduce errore di
approssimazione, ma riduce porte e profondità. Per $N=21$, il troncamento nelle
trasformate interne agli addizionatori porta le porte 2Q da 49\,661 a 23\,178;
nella stima di fase la variante troncata supera quella completa in tutti i nove
confronti eseguiti. È un risultato esplorativo, ripreso nella
Sezione~\ref{sec:qft_topologia}; confronto con le proposte che riducono i qubit
logici e relativi costi sono nella Sezione~\ref{app:aqft_estesa}.

\section{Dall'esito della prima fase alla struttura del lavoro}

Il filo sperimentale procede quindi da un controllo NISQ appaiato, che separa
effetto TOP-K ed effetto ML, alla protezione durante il calcolo: codice a
ripetizione, Steane, surface code e decodifica. La sensibilità di Shor al proxy
per gate resta distinta dai tassi logici per ciclo; l'\ac{AQFT} chiude il
percorso agendo direttamente sul costo circuitale.

\end{onehalfspace}
